\documentclass[a4paper,10pt]{report}\input{../0.Préambule/Préambule_cours_prof}\begin{document}

\newpage
\section*{Les translations}
\addcontentsline{toc}{section}{Les translations}

\begin{exer1}
Le pavage ci-dessous est réalisé avec $30$ pièces identiques.

\noindent \begin{minipage}{9cm}
\includegraphics[scale=0.45]{../40.Image/19. Exo-Translations_1.png}
\end{minipage}
\hspace{5mm}
\begin{minipage}{7cm}
Observe le pavage puis réponds aux questions :
\begin{enumerate}[font=\color{exer1}]
\item Dans la translation qui transforme $A$ en $B$ :
\begin{enumerate}[font=\color{exer1}]
\item[•] Quelle est l'image de la pièce $13$ ?
\item[•] Quelle est l'image de la pièce $6$ ?
\item[•] Quelle est l'image de la pièce $15$ ?
\item[•] Quelle est l'image de la pièce $1$ ?
\end{enumerate}
\item Dans la translation qui transforme $A$ en $B$ :
\begin{enumerate}[font=\color{exer1}]
\item[•] Quelle est l'image de la pièce $25$ ?
\item[•] Quelle est l'image de la pièce $18$ ?
\item[•] Quelle est l'image de la pièce $23$ ?
\item[•] Quelle est l'image de la pièce $20$ ?
\end{enumerate}
\item Quelle remarque peux-tu faire au sujet de ces deux translations ?
\item Dans la translation qui transforme $C$ en $F$ :
\begin{enumerate}[font=\color{exer1}]
\item[•] Quelle est l'image du point $D$ ?
\item[•] Place $P$, image de $N$
\item[•] Place le point $E$ qui a pour image $N$
\end{enumerate}
\end{enumerate}
\end{minipage}
\end{exer1}

\begin{exer1}

\medskip
\begin{enumerate}[font=\color{exer1}]
\item Colorie en rouge le motif qui permet \\d'obtenir cette frise en effectuant \\ la translation qui transforme $C$ en $D$.

\vspace{-2.5cm}
\begin{flushright}
\includegraphics[scale=0.4]{../40.Image/19. Exo-Translations_2.png}
\end{flushright}

\item Représente par une flèche la translation effectué \\de la frise pour obtenir le pavage ci-dessous, puis \\ colorie ce pavage en alternant deux couleurs de ton \\ choix pour chaque motif obtenu par translation.

\vspace{-2.5cm}
\begin{flushright}
\includegraphics[scale=0.4]{../40.Image/19. Exo-Translations_3.png}
\end{flushright}
\end{enumerate}
\end{exer1}

\noindent \begin{minipage}{8.3cm}
\begin{exer1}
Dans chaque cas, trace en rouge l'image du triangle par la translation qui transforme $A$ en $B$ et en vert l'image du triangle par la translation qui transforme $C$ en $D$.

\vspace{8mm}
\includegraphics[scale=0.4]{../40.Image/19. Exo-Translations_4.png}
\end{exer1}
\end{minipage}
\hspace{2mm}
\vline
\hspace{2mm}
\begin{minipage}{8.5cm}
\begin{exer1}

\medskip
\includegraphics[scale=0.4]{../40.Image/19. Exo-Translations_5.png}

\begin{enumerate}[font=\color{exer1}]
\item Dessine en rouge l'image du bateau par la translation qui transforme $K$ en $P$.
\item Dessine en vert l'image du bateau par la translation qui transforme $L$ en $P$.
\end{enumerate}
\end{exer1}
\end{minipage}

\begin{exer1}

\vspace{-5mm}
\begin{center}
\includegraphics[scale=0.4]{../40.Image/19. Exo-Translations_6.png}
\end{center}

\begin{enumerate}[font=\color{exer1}]
\item Trace $A'B'C'D'E'F'$, l'image du polygone $ABCDEF$ par la translation qui transforme $B$ en $A$.
\item Trace $A''B''C''D''E''F''$, l'image du polygone $ABCDEF$ par la translation qui transforme $A$ en $E$.
\item Reproduis ces translations avec les deux polygones obtenus et poursuis ainsi le dessin.
\item Comment s'appelle le dessin obtenu ?
\end{enumerate}
\end{exer1}

\noindent \begin{minipage}{7.8cm}
\begin{exer1}
Construis l'image du quadrilatère :
\begin{enumerate}[font=\color{exer1}]
\item par translation qui transforme $D$ en $D'$
\item par la translation qui transforme $B$ en $D$.
\end{enumerate}

\vspace{2.2cm}
\includegraphics[scale=0.3]{../40.Image/19. Exo-Translations_7.png}
\vspace{2.2cm}.
\end{exer1}
\end{minipage}
\hspace{2mm}
\vline
\hspace{2mm}
\begin{minipage}{8.5cm}
\begin{exer1}
Construis une frise en appliquant au motif rouge une symétrie axiale d'axe $(AE)$ puis en appliquant à ces deux motifs la translation qui transforme $A$ en $E$.

\vspace{3mm}
\includegraphics[scale=0.4]{../40.Image/19. Exo-Translations_8.png}
\end{exer1}

\begin{exer2}
\begin{enumerate}[font=\color{exer2}]
\item Trace ne rouge l'image de cette figure par la translation qui transforme $C$ en $H$.
\item Trace en vert l'image de cette figure par la symétrie de centre $I$
\item Trace en bleu l'image de la figure verte pas la symétrie de centre $J$.
\item Par quelle transformation passe-t-on de la figure noire à la figure bleue ?
\end{enumerate}
\end{exer2}
\end{minipage}

\vspace{1mm}\hspace{3cm}\includegraphics[scale=0.35]{../40.Image/19. Exo-Translations_9.png}

\end{document}